\documentclass{article}
\usepackage{amsmath, amssymb, amsthm}
\usepackage{hyperref}
\usepackage{geometry}
\geometry{margin=1in}

\title{Erd\H{o}s Problem \#1052}
\date{Last updated: 2025-09-28}
\author{Source: erdosproblems.com}

\begin{document}
\maketitle

\noindent\textbf{Status:} OPEN \quad \textbf{Prize: \$10}

\noindent\textbf{Tags:} number theory

\medskip
\noindent\textbf{Note:} unitary perfect numbers

\medskip

\noindent\textbf{Problem Statement:}

A unitary divisor of $n$ is $d\mid n$ such that $(d,n/d)=1$. A number $n\geq 1$ is a unitary perfect number if it is the sum of its unitary divisors (aside from $n$ itself). Are there only finite many unitary perfect numbers?




Guy \cite{Gu04} reports that Carlitz, Erd\H{o}s, and Subbarao offer \$10 for settling this question, and that Subbarao offers 10 cents for each new example. There are no odd unitary perfect numbers. There are five known unitary perfect numbers ( A002827 in the OEIS):\[6, 60, 90, 87360, 146361946186458562560000.\]This is problem B3 in Guy's collection \cite{Gu04}.




References

[Gu04] Guy, Richard K., Unsolved problems in number theory . (2004), xviii+437.

\noindent\textbf{Related OEIS sequences:} A002827


\bigskip
\noindent\small{Source: \url{https://www.erdosproblems.com/1052}}
\end{document}
